\chapter{Магнітні поверхні, $q$ і кут прямих силових ліній}\label{ch:geo3}

\section{Поверхня як контур $\psi$}\label{sec:geo3-surf}

Магнітна поверхня в перерізі --- замкнена лінія рівня $\psiN=s$, що охоплює вісь. Лінія рівня на дискретній сітці --- це стандартна задача «marching squares»; проєкт бере її з \texttt{scikit-image} і відсіює все зайве:

\code{viz/src/tokviz/equilibrium.py}{250}{280}{(поверхня $\psiN=s$)}

\noindent Рядок 260: усі контури рівня $s$. На одному рівні їх може бути кілька: крім поверхні навколо осі, трапляються контури в приватній області під X-точкою або біля котушок. Рядки 263--265: дробові індекси сітки переводяться в метри. Рядки 266--273 відкидають незамкнені контури (вони виходять на межу сітки), короткі ($<16$ точок) і ті, що не охоплюють вісь (\texttt{\_encircles}: кут навколо осі набігає більш ніж на $1{,}5\pi$). Рядки 278--279 орієнтують контур проти годинникової стрілки за знаком площі. Орієнтація потрібна, щоб полоїдальний кут зростав у передбачуваному напрямку.

Під X-точкою (розд.~\ref{ch:geo5}) замкненого контуру навколо осі при $\psiN\ge1$ немає. Тому \texttt{surface(1.0)} на діверторній рівновазі ненадійна, і проєкт працює з $\psiN\le0{,}98$--$0{,}99$. Звідси й звичай характеризувати край плазми величиною $\qnf=q(\psiN=0{,}95)$, а не $q$ на самій межі: на сепаратрисі $q\to\infty$ (п.~\ref{sec:geo5-x}).

\section{Запас стійкості $q$}\label{sec:geo3-q}

$q$ --- скільки разів силова лінія обходить тор уздовж $\tor$ за один полоїдальний оберт. Уздовж лінії $R\,\dd\tor/\dd\ell=\Btor/\Bpol$, де $\dd\ell$ --- елемент полоїдальної дуги. За один полоїдальний оберт з~\eqref{eq:geo2-comp}
\begin{equation}\label{eq:geo3-q}
\Delta\tor=\oint\frac{\Btor}{R\,\Bpol}\,\dd\ell=F\oint\frac{\dd\ell}{R\,|\grad\psi|},\qquad q=\frac{\Delta\tor}{2\pi}=\frac{F}{2\pi}\oint\frac{\dd\ell}{R\,|\grad\psi|}.
\end{equation}
Це та сама формула, що й $q=\frac1{2\pi}\oint F|\mathcal{J}|R^{-2}\dd\theta$ у методичці~\cite[розд.~4, (4.1)]{Manual}: для кута, що вимірює довжину дуги, якобіан $|\mathcal{J}|=R\,\dd\ell/(|\grad\psi|\,\dd\theta)$. Інтеграл береться чисельно, сумою по відрізках контуру:

\code{viz/src/tokviz/equilibrium.py}{283}{306}{(контурний інтеграл і $q$)}

\noindent Рядки 290--292: довжини відрізків $\dd\ell$ замкненого багатокутника. Рядки 293--297: вага $\dd\ell/(R|\grad\psi|)$ у середині кожного відрізка; рядок 296 захищає від ділення на нуль біля X-точки, де $|\grad\psi|\to0$. Рядок 298: накопичена сума \texttt{cum} --- вона ще знадобиться для $\thetastar$. Рядок 306: $q$ за~\eqref{eq:geo3-q}.

Формула~\eqref{eq:geo3-q} лінійна за $F$, тому $F$ відновлюють з відомого $\qnf$ одним діленням (п.~\ref{sec:geo2-F}):

\code{viz/src/tokviz/equilibrium.py}{308}{322}{(калібрування $F$ за $\qnf$)}

\begin{center}\small
\begin{tabular}{@{}lccccccc@{}}\toprule
$\psiN$ & 0,10 & 0,30 & 0,50 & 0,70 & 0,90 & 0,95 & 0,98\\\midrule
$q$ (DIII-D \#203702, кадр 75) & 0,759 & 0,956 & 1,262 & 1,791 & 2,999 & 3,699$^{*}$ & 4,497\\\bottomrule
\end{tabular}\\[2pt]
{\footnotesize $^{*}$ --- за побудовою (калібрування $F$). Запуск \texttt{geom\_tour.py}.}
\end{center}

Профіль типовий для діверторної плазми: $q$ повільно зростає в ядрі й круто --- біля межі, де поверхні проходять поруч з X-точкою і $|\grad\psi|$ малий. У центрі $q<1$: поверхня $q=1$ лежить біля $\psiN\approx0{,}33$ (саме там у ранній версії візуалізації з'явився фальшивий острів, VE6 нижче).

\begin{established}{VE4, VE16}
Два незалежні способи дають те саме $q$: контурний інтеграл~\eqref{eq:geo3-q} і пряме трасування силової лінії (тороїдальні оберти на полоїдальні) збігаються до 0,001--0,16\,\% на $\psiN=0{,}3\ldots0{,}95$ (VE4). Число обертання $1/q$, виміряне трасуванням, відтворює рівноважне до 0,03\,\% (VE16)~\cite{RegViz}.
\end{established}

\section{Кут прямих силових ліній $\thetastar$}\label{sec:geo3-thetastar}

У геометричному куті $\theta=\arctan[(Z-\Zax)/(R-\Rax)]$ силова лінія на поверхні --- крива: на внутрішньому боці тора вона йде полоїдально повільніше, ніж на зовнішньому. Кут прямих силових ліній $\thetastar$ вибирають так, щоб лінія просувалася рівномірно:
\begin{equation}\label{eq:geo3-thetastar}
\frac{\dd\thetastar}{\dd\tor}=\frac1q\quad\Longrightarrow\quad
\thetastar(\ell)=2\pi\,\frac{\int_0^{\ell}\dd\ell'/(R|\grad\psi|)}{\oint\dd\ell'/(R|\grad\psi|)} .
\end{equation}
Друга рівність випливає з першої: з~\eqref{eq:geo3-q} на відрізку $\dd\ell$ лінія проходить $\dd\tor=F\,\dd\ell/(R|\grad\psi|)$, а $\dd\thetastar=\dd\tor/q$. Отже, $\thetastar$ --- це просто нормована накопичена сума \texttt{cum} з контурного інтеграла. $F$ скорочується, тому $\thetastar$ від $F$ не залежить.

\begin{figure}[!tb]\centering
\includegraphics[width=0.86\textwidth]{geo_thetastar.jpg}
\caption{Лінії сталого кута (червоні, крок $\pi/8$) на поверхнях $\psiN=0{,}1\ldots0{,}9$ DIII-D \#203702, кадр 75. Ліворуч --- геометричний кут $\theta$ (промені з осі), праворуч --- кут прямих силових ліній $\thetastar$, відлічений від зовнішнього екватора. Лінії $\thetastar$ густі на внутрішньому боці: там силова лінія проводить більше тороїдального «часу» на одиницю полоїдальної дуги. Власний розрахунок, код: \texttt{geometry/code/geom\_tour.py}.}\label{fig:geo3-thetastar}
\end{figure}

Навіщо стільки клопоту? Збурення вигляду $\cos(m\thetastar-n\tor)$ на поверхні, де $q=m/n$, має сталу фазу вздовж силової лінії і тому \emph{резонує}: саме там виростає магнітний острів $m/n$. У геометричному куті така гармоніка «розмазується» по кількох $m$ і дає острови не там (VR1 нижче). На рис.~\ref{fig:geo3-thetastar} видно, наскільки кути різні: найбільше відхилення $|\thetastar-\theta|$ становить $15^\circ$ на $\psiN=0{,}2$, $25^\circ$ на $0{,}5$ і $48^\circ$ на $0{,}9$ (з наближенням до X-точки воно зростає).

\code{viz/src/tokviz/equilibrium.py}{363}{395}{(таблиця $\Delta=\thetastar-\theta$)}

\noindent Код табулює не сам $\thetastar$, а різницю $\Delta=\thetastar-\theta$ на регулярній сітці $(\psiN,\theta)$. Рядок 370: $\thetastar$ з накопиченої суми~\eqref{eq:geo3-thetastar}. Рядки 371--372: геометричний кут кожної точки контуру, «розгорнутий» без стрибків $2\pi$. Рядки 376--380 пропускають поверхні, де $\theta$ уздовж контуру не монотонний (біля X-точки поверхня перестає бути «зіркоподібною» відносно осі). Рядки 384--390: $\Delta$ однозначна і $2\pi$-періодична, бо обидва кути набігають рівно $2\pi$ за оберт; її інтерполюють на рівномірну сітку $\theta$. Похідні $\thetastar$ по $R$ і $Z$ далі беруть за правилом ланцюга (клас \texttt{ThetaStarField}, рядки 410--481 файлу).

\begin{established}{VE6, VR1, VR2}
Перша реалізація растеризувала $\cos(m\thetastar)$ на сітку $65\times65$. Це дало 2,4\,\% паразитної гармоніки $m=1$ і фальшивий ланцюг островів на $q=1$ ($\psiN=0{,}335$), хоча збурення містило лише $m=2$ і $3$. Аналітичне обчислення через табульовану $\Delta$ зменшило паразитну $m=1$ до 0,4\,\%, і фальшивий острів зник (VE6). Звідси два спростування: «геометричного кута достатньо для розміщення островів» (VR1) і «паразитні гармоніки з інтерполяції можна ігнорувати» (VR2)~\cite{RegViz}.
\end{established}

\begin{newfinding}[початок відліку $\thetastar$]
$\thetastar$ визначено з точністю до сталої на кожній поверхні. У \texttt{theta\_star\_geometry} ця стала --- кут першої точки контуру, яку повертає \texttt{find\_contours} (рядок 370: \texttt{cum} починається з нуля там, де почався контур). Ця точка на сусідніх поверхнях різна. Запуск \texttt{geom\_tour.py} на кадрі 75 дає: $\Delta(\theta=0)$ по 200 поверхнях таблиці лежить у межах $-79{,}7^\circ\ldots-36{,}6^\circ$, а між сусідніми поверхнями ($\delta\psiN\approx0{,}005$) стрибає до $34{,}6^\circ$.

Для $\thetastar$ на одній поверхні це неважливо, але \texttt{ThetaStarField.value\_and\_grad} бере ще й $\pa\Delta/\pa\psiN$ (рядки 479--480 файлу), і стрибки потрапляють у градієнт $\thetastar$, а через нього --- у поле збурення (\texttt{perturbation.py}, рядки 156--157). Середньоквадратичне $|\grad\thetastar|$ уздовж поверхні, рад/м:
\begin{center}\small
\begin{tabular}{@{}lccc@{}}\toprule
$\psiN$ & 0,3 & 0,6 & 0,9\\\midrule
як у коді & 38,0 & 3,9 & 3,03\\
з відліком від зовнішнього екватора & 3,5 & 2,5 & 2,89\\\bottomrule
\end{tabular}
\end{center}
На $\psiN=0{,}3$ градієнт завищено в 11 разів, на $\psiN=0{,}9$, де лежать робочі резонанси проєкту ($q=3$, $10/3$, $11/3$), --- на 5\,\%. Це узгоджується з тим, що виміряні ширини островів збігаються з формулою (VE18). \emph{Вплив на реєстр VE/VR не перевірено.} Можливе виправлення в один рядок: після рядка 390 віднімати від кожного рядка таблиці $\Delta$ його значення при $\theta=0$. Так зроблено в \texttt{geom\_tour.py}, рядки 104--109, і на рис.~\ref{fig:geo3-thetastar}.
\end{newfinding}

\section{Від перерізу до тора}\label{sec:geo3-3d}

Осесиметрія робить 3D-геометрію тривіальною: будь-який контур $(R,Z)$ --- межа плазми, поверхня $\psiN$, профіль котушки --- обертають навколо осі $Z$. Точка $(R,Z)$ на куті $\tor$ переходить у $(x,y,z)=(R\cos\tor,R\sin\tor,Z)$:

\code{viz/src/tokviz/surfaces.py}{38}{68}{(поверхня обертання)}

\noindent Рядки 52--55 --- саме ця формула для всіх точок профілю й усіх кутів одразу. Рядки 57--67 збирають чотирикутні грані між сусідніми кільцями; \texttt{\% n\_tor} і \texttt{\% n\_pol} замикають сітку по обох кутах. Параметр \texttt{scale} множить $R$ і $Z$ однаково. Для рендера проєкт збільшує DIII-D у 3 рази (\texttt{GEOMETRY\_SCALE} у \texttt{config.py}): аспектне відношення при цьому не змінюється, а отже й форма поверхонь та розташування островів у нормованих координатах.

\begin{practicum}[$q$ і $\thetastar$]
(1)~Порахуйте $q$ на $\psiN=0{,}5$ контурним інтегралом при \texttt{surface()} на сітці, згущеній удвічі (інтерполюйте $\psi$ сплайном на $129\times129$ і збудуйте нову \texttt{Equilibrium}). Наскільки зміниться $q$? (2)~Знайдіть $\psiN$ поверхні $q=1$ бісекцією по \texttt{q\_at}. (3)~Повторіть таблицю «Нового спостереження» для кадрів 10 і 150 того самого розряду. Чи змінюються стрибки початку відліку?
\end{practicum}
